\chapter{Diagrammes de séquence}
\label{chap:sequences}

\section{Ouverture d'une image}

\begin{figure}[htbp]
\centering
\begin{tikzpicture}[x=1cm,y=1cm,font=\small]
  \foreach \x/\name in {0/Utilisateur,2.1/QML,4.2/ImageEditor,6.5/Contrôleur,8.4/Import,10.7/Sidecar} {
    \node[align=center] at (\x,0) {\name};
    \draw[densely dashed,gray] (\x,-0.4) -- (\x,-8.0);
  }
  \draw[message] (0,-1) -- node[above,font=\footnotesize]{ouvre une image} (2.1,-1);
  \draw[message] (2.1,-2) -- node[above,font=\footnotesize]{openImage()} (4.2,-2);
  \draw[message] (4.2,-2.65) -- node[above,font=\footnotesize]{lire recette RAW/WB .ipr} (10.7,-2.65);
  \draw[message] (4.2,-3.25) -- node[above,font=\footnotesize]{acquire()} (6.5,-3.25);
  \draw[message] (6.5,-3.65) -- node[above,font=\footnotesize]{admission} (4.2,-3.65);
  \draw[message] (4.2,-4.45) -- node[above,font=\footnotesize]{ImageImport::load()} (8.4,-4.45);
  \draw[message] (8.4,-5.1) -- node[above,font=\footnotesize]{WorkingImage} (4.2,-5.1);
  \node[align=left,anchor=west,font=\scriptsize] at (4.3,-5.8) {retient source active + proxy};
  \draw[message] (4.2,-6.5) -- node[above,font=\footnotesize]{lire .ipr si valide} (10.7,-6.5);
  \draw[message] (10.7,-7.15) -- node[above,font=\footnotesize]{Settings + métadonnées} (4.2,-7.15);
  \draw[message] (4.2,-7.75) -- node[above,font=\footnotesize]{aperçu prêt} (2.1,-7.75);
\end{tikzpicture}
\caption{Chemin cache miss admis, simplifié. L'attente mémoire ne bloque pas QML. La réservation est restituée par le worker ; les résultats périmés sont abandonnés. L'image quittée peut rejoindre le LRU, pas systématiquement la nouvelle image active.}
\end{figure}

\section{Réglage et application par lot}

\begin{figure}[htbp]
\centering
\begin{tikzpicture}[x=1cm,y=0.8cm,font=\small]
  \foreach \x/\name in {0/Utilisateur,2.2/Setting.qml,4.8/ImageEditor,7.4/Pipeline,10.0/Sidecars} {
    \node[align=center] at (\x,0) {\name};
    \draw[densely dashed,gray] (\x,-0.4) -- (\x,-7.2);
  }
  \draw[message] (0,-1) -- node[above,font=\footnotesize]{modifie un réglage} (2.2,-1);
  \draw[message] (2.2,-1.8) -- node[above,font=\footnotesize]{setter} (4.8,-1.8);
  \draw[message] (4.8,-2.6) -- node[above,font=\footnotesize]{rendu différé} (7.4,-2.6);
  \draw[message] (7.4,-3.25) -- node[above,font=\footnotesize]{image rendue} (4.8,-3.25);
  \draw[message] (4.8,-3.95) -- node[above,font=\footnotesize]{aperçu} (2.2,-3.95);
  \draw[message] (4.8,-4.65) -- node[above,font=\footnotesize]{écriture atomique, 450 ms} (10.0,-4.65);
  \draw[message] (0,-5.45) -- node[above,font=\footnotesize]{appliquer à N images} (2.2,-5.45);
  \draw[message] (2.2,-6.15) -- node[above,font=\footnotesize]{applyAdjustmentsToSelection()} (4.8,-6.15);
  \draw[message] (4.8,-6.85) -- node[above,font=\footnotesize]{un sidecar par cible} (10.0,-6.85);
\end{tikzpicture}
\caption{Cycle de réglage non destructif et propagation volontaire à une sélection.}
\end{figure}

\section{Changement de profil et normalisation}

\begin{figure}[htbp]
\centering
\begin{tikzpicture}[x=1cm,y=0.8cm,font=\small]
  \foreach \x/\name in {0/Module QML,3.1/ImageEditor,6.5/Worker RAW,10.2/Profils et manifeste} {
    \node[align=center] at (\x,0) {\name};
    \draw[densely dashed,gray] (\x,-0.4) -- (\x,-7.3);
  }
  \draw[message] (0,-1) -- node[above]{setCameraProfile()} (3.1,-1);
  \draw[message] (3.1,-2) -- node[above]{tampon caméra ou décodage} (6.5,-2);
  \draw[message] (6.5,-3) -- node[above]{applyProfile() ou resolve()} (10.2,-3);
  \draw[message] (10.2,-4) -- node[above]{profil / calibration ou refus} (6.5,-4);
  \draw[message] (6.5,-5) -- node[above]{image + statut réel} (3.1,-5);
  \draw[message] (3.1,-6) -- node[above]{retouches conservées} (0,-6);
  \draw[message] (0,-7) -- node[above]{décocher : clé non calibrée} (3.1,-7);
\end{tikzpicture}
\caption{Le tampon pré-DCP évite LibRaw quand WB/calibration/optique le permettent. Le rendu visible reste conservé, la recette v6 est sauvegardée après succès. La case seule ne garantit pas une calibration.}
\end{figure}
